\documentclass[11pt,a4paper]{article}

% ---------- Encoding & Sprache ----------
\usepackage[utf8]{inputenc}
\usepackage[T1]{fontenc}
\usepackage[ngerman]{babel}
\usepackage{lmodern}

% ---------- Layout ----------
\usepackage[margin=2cm]{geometry}
\usepackage{setspace}
\usepackage{enumitem}
\usepackage{booktabs}
\usepackage{hyperref}
\setlist{nosep, itemsep=1pt, topsep=2pt}

\setlength{\parindent}{0pt}
\setlength{\parskip}{4pt}

\title{\vspace{-1.5cm}\Large\textbf{Portfolioteil 3 -- Finalisierungsphase}\\[2pt]
\large Funktionalität und Metadaten der Datenbank\\[2pt]
\normalsize DLBDSPBDM01\_D -- Data-Mart-Erstellung in SQL}
\author{Tizian Senger \\ Matrikelnummer: IU14143428}
\date{\today}

\begin{document}
\maketitle
\vspace{-1cm}

% =========================================================
\section*{1. Funktionalität des Datenbankmanagementsystems}
% =========================================================
% TODO: Beschreibe hier, was die finale Datenbank kann (Buch anbieten, ausleihen,
% Umkreissuche, Bewertungen, Benachrichtigungen etc.) und welche Verbesserungen
% gegenueber Phase 2 aufgrund des Tutor:innen-Feedbacks vorgenommen wurden.

% =========================================================
\section*{2. Metadaten der Datenbank}
% =========================================================
% TODO: Werte nach finaler Befuellung eintragen.
\begin{center}
\begin{tabular}{@{}ll@{}}
\toprule
\textbf{Metadatum} & \textbf{Wert} \\
\midrule
Anzahl Tabellen & TODO \\
Anzahl Datensätze gesamt & TODO \\
Größe der Datenbank & TODO \\
Verwendetes DBMS & TODO \\
\bottomrule
\end{tabular}
\end{center}

% =========================================================
\section*{3. GitHub-Repository}
% =========================================================
% TODO: Link zum GitHub-Repository mit Code, Berichten und README eintragen.
Repository: \url{https://github.com/TODO/TODO}

\end{document}
